% Complete categorical chains; hypotheses and existence policy are stated in the body.
\begin{annexitem}{annex:p13-boger-conversions}{thm:p13-boger-conversions}{Conversion and the existence dependency}
\begin{proof}
The terms are homogeneous: the same subject matter can be named by either
term, rather than one term naming a sentence and the other an individual.
We use the ordinary predication meanings stated in the body and in Boger's
pp.~165,168--169,180--183 account
\parencite{Boger2004UnderlyingLogic}.

\emph{CR1.} Suppose \(AeB\) and temporarily suppose its proposed
converse is false. The latter gives \(BiA\): some thing named by \(A\)
also has the property expressed by \(B\). Select that thing, as in
Aristotle's setting-out argument. It is a \(B\) to which \(A\) belongs,
so \(AiB\) holds, contradicting \(AeB\). The assumed failure of
\(BeA\) is impossible; \(BeA\) follows. No identity \(C=A\) between
entire term extensions was used. The selection is a shared positive
case, the feature for which the source's note~70 qualifies its notation.

\emph{CR2.} Suppose \(AaB\). For the source's indirect route, suppose
\(BeA\), contradictory to the intended \(BiA\). CR1 gives \(AeB\).
The actual universal affirmation provides existent \(B\), so select a
\(B\). The actual premise attributes \(A\) to it; \(AeB\) denies that
attribution to every \(B\). The two universals are incompatible for
this existent subject. Discharge \(BeA\), obtaining \(BiA\).
Equivalently the selected existent \(B\) is itself the shared positive
case. Without existent \(B\), a vacuous universal would not justify
this inference.

\emph{CR3.} Suppose \(AiB\) and temporarily suppose \(BeA\).
CR1 gives \(AeB\), contradictory to \(AiB\). Discharge \(BeA\),
obtaining \(BiA\). Here the positive particular already supplies the
existent common case, so no additional universal-subject premise is needed.

For both failures use the source's interpreted situation in which every
man is an animal, some animals are not men, and men exist. With \(A\)
animal and \(B\) man, \(AaB\) is true while \(BaA\) is false.
With \(A\) man and \(B\) animal, \(AoB\) is true while \(BoA\)
is false, because every man is an animal. These are genuine counterinstances
to the two proposed simple conversions, not failures to find a proof.
\end{proof}
\end{annexitem}

\begin{annexitem}{annex:p13-boger-ten-completions}{thm:p13-boger-ten-completions}{The ten completion deductions}
\begin{proof}
Every letter below retains its predicate-first role. CR1, CR2, and CR3 mean
the e, a-per-accidens, and i conversions. Each first-figure application uses
the earlier formulas in precisely the indicated predicate/middle/subject
positions. The original source is table~19, p.~187
\parencite{Boger2004UnderlyingLogic}; we remove unnecessary repetitions but
do not replace the categorical steps by first-order quantifier reasoning.

\emph{Cesare.} Premises \(MeN,MaX\). CR1 on \(MeN\) gives \(NeM\).
Celarent on \(NeM,MaX\) gives \(NeX\), the required conclusion.

\emph{Camestres.} Premises \(MaN,MeX\). CR1 gives \(XeM\).
Celarent on \(XeM,MaN\) gives \(XeN\). CR1 gives \(NeX\).

\emph{Festino.} Premises \(MeN,MiX\). CR1 gives \(NeM\).
Ferio on \(NeM,MiX\) gives \(NoX\).

\emph{Baroco.} Premises \(MaN,MoX\). Open the single temporary assumption
\(NaX\), contradictory to the intended \(NoX\). Barbara on \(MaN,NaX\)
gives \(MaX\). This and the actual premise \(MoX\) are contradictories.
Close that subdeduction and discharge \(NaX\), obtaining \(NoX\).

\emph{Darapti.} Premises \(PaS,RaS\). CR2 on \(RaS\) gives \(SiR\),
using the existent \(S\) carried by that universal affirmation.
Darii on \(PaS,SiR\) gives \(PiR\).

\emph{Felapton.} Premises \(PeS,RaS\). CR2 on \(RaS\) gives \(SiR\).
Ferio on \(PeS,SiR\) gives \(PoR\).

\emph{Disamis.} Premises \(PiS,RaS\). CR3 on \(PiS\) gives \(SiP\).
Darii on \(RaS,SiP\) gives \(RiP\). CR3 gives \(PiR\), restoring
the required orientation of the extremes.

\emph{Datisi.} Premises \(PaS,RiS\). CR3 on \(RiS\) gives \(SiR\).
Darii on \(PaS,SiR\) gives \(PiR\).

\emph{Bocardo.} Premises \(PoS,RaS\). Temporarily assume \(PaR\),
the contradictory of the intended \(PoR\). Barbara on \(PaR,RaS\)
gives \(PaS\), contradictory to the actual premise \(PoS\).
Discharge \(PaR\), obtaining \(PoR\).

\emph{Ferison.} Premises \(PeS,RiS\). CR3 on \(RiS\) gives \(SiR\).
Ferio on \(PeS,SiR\) gives \(PoR\).

These exhaust SR5--14. The only temporary assumptions occur in Baroco
and Bocardo, and both have been discharged. No derived conclusion of one
case has been imported as an unproved premise into another case.
\end{proof}
\end{annexitem}

\begin{annexitem}{annex:p13-boger-contrast-four}{thm:p13-boger-contrast-four}{The four conclusions excluded by two contrasts}
\begin{proof}
Use the substitution with true premises and \(PeS\) to exclude \(PiS\),
its contradictory. Because the subject exists, \(PeS\) is incompatible
with \(PaS\), so that same substitution also excludes \(PaS\).
Use the other substitution, with true premises and \(PaS\), to exclude
\(PoS\), its contradictory. Existence makes \(PaS\) incompatible with
\(PeS\), so the latter substitution also excludes \(PeS\).
Every candidate is therefore false in at least one admissible substitution
that makes both premises true. None follows necessarily from the pattern.
Both substitutions are used; truth of one displayed conclusion alone
does not supply the assertion for all four candidates.
\end{proof}
\end{annexitem}

\begin{annexitem}{annex:p13-boger-weakening-inconcludence}{thm:p13-boger-weakening-inconcludence}{Keeping each counterinstance after weakening}
\begin{proof}
Fix one candidate conclusion and its counterinstance for
\(\Gamma_{\mathrm{s}}\). Its conclusion is false and both stronger premises
are true. By the stipulated implication between the premise conditions,
both weaker premises are also true in that very substitution. It is
therefore a counterinstance for \(\Gamma_{\mathrm{w}}\), with no change
of conclusion or term orientation. Repeat for the four candidates.
For the stated categorical weakening, existent subject matter and
indeterminate some make \(MeS\) imply \(MoS\): if no \(S\) has \(M\)
and at least one \(S\) exists, at least one \(S\) lacks \(M\).
The unchanged other premise completes the implication between the pairs.
This last implication fails for some-but-not-all because a universal
denial supplies no positive \(S\)-case.
\end{proof}
\end{annexitem}

\begin{annexitem}{annex:p13-boger-reduced-basis}{thm:p13-boger-reduced-basis}{Twelve reductions and finite expansion}
\begin{proof}
Barbara and Celarent are the retained basis. The following twelve deductions
use only them, admitted conversions, and discharged assumptions. They
reconstruct the routes discussed on pp.~209--214 and do not certify
unread clipped entries as literal transcriptions
\parencite{Boger2004UnderlyingLogic}.

\emph{Darii.} From \(AaB,BiC\), assume \(AeC\). CR1 gives \(CeA\).
Celarent with \(CeA,AaB\) gives \(CeB\); CR1 gives \(BeC\).
That contradicts \(BiC\). Discharge \(AeC\) to obtain \(AiC\).

\emph{Ferio.} From \(AeB,BiC\), assume \(AaC\). CR1 gives \(BeA\).
Celarent with \(BeA,AaC\) gives \(BeC\), contradictory to \(BiC\).
Discharge \(AaC\) to obtain \(AoC\).

\emph{Cesare.} From \(MeN,MaX\), CR1 gives \(NeM\), and Celarent
with \(NeM,MaX\) gives \(NeX\).

\emph{Camestres.} From \(MaN,MeX\), CR1 gives \(XeM\); Celarent
with \(XeM,MaN\) gives \(XeN\); CR1 gives \(NeX\).

\emph{Festino.} From \(MeN,MiX\), assume \(NaX\).
Celarent with \(MeN,NaX\) gives \(MeX\), contradictory to \(MiX\).
Discharge \(NaX\) to obtain \(NoX\). No Ferio application is hidden here.

\emph{Baroco.} From \(MaN,MoX\), assume \(NaX\).
Barbara with \(MaN,NaX\) gives \(MaX\), contradictory to \(MoX\).
Discharge \(NaX\) to obtain \(NoX\).

\emph{Darapti.} From \(PaS,RaS\), assume \(PeR\).
CR1 gives \(ReP\). Celarent with \(ReP,PaS\) gives \(ReS\).
The universal affirmation \(RaS\) provides existent \(S\), so
\(ReS,RaS\) are incompatible. Discharge \(PeR\) to obtain \(PiR\).

\emph{Felapton.} From \(PeS,RaS\), assume \(PaR\).
CR1 gives \(SeP\). Celarent with \(SeP,PaR\) gives \(SeR\);
CR1 gives \(ReS\). This conflicts with \(RaS\) for its existent
subject \(S\). Discharge \(PaR\) to obtain \(PoR\).

\emph{Disamis.} From \(PiS,RaS\), assume \(PeR\).
Celarent with \(PeR,RaS\) gives \(PeS\), contradictory to \(PiS\).
Discharge \(PeR\) to obtain \(PiR\).

\emph{Datisi.} From \(PaS,RiS\), assume \(PeR\).
CR1 gives \(ReP\); Celarent with \(ReP,PaS\) gives \(ReS\),
contradictory to \(RiS\). Discharge \(PeR\) to obtain \(PiR\).

\emph{Bocardo.} From \(PoS,RaS\), assume \(PaR\).
Barbara with \(PaR,RaS\) gives \(PaS\), contradictory to \(PoS\).
Discharge \(PaR\) to obtain \(PoR\).

\emph{Ferison.} From \(PeS,RiS\), assume \(PaR\).
CR1 gives \(SeP\). Celarent with \(SeP,PaR\) gives \(SeR\);
CR1 gives \(ReS\), contradictory to \(RiS\).
Discharge \(PaR\) to obtain \(PoR\).

For the final expansion claim, process a finite full-system deduction from
its first line to its last. Keep premise and repetition lines. Keep conversion,
Barbara, and Celarent applications. Replace each application of one of
the other twelve rules by its just-proved finite chain with the same letters
substituted uniformly. All required premises are earlier available lines.
Each replacement has the same conclusion and has discharged all its own
temporary assumptions. For a line inside an existing indirect subdeduction,
perform the replacement inside that subdeduction, without discharging
its enclosing assumption. Thus earlier availability and assumption scope
are preserved at every step. There are finitely many original lines and
each replacement is finite; the expanded deduction is finite. Its last
line is unchanged and its undischarged premises are precisely those of
the original. This proves sufficiency of the reduced basis, not an
independence or semantic completeness theorem.
\end{proof}
\end{annexitem}

\begin{annexitem}{annex:p13-boger-analysis-links}{thm:p13-boger-analysis-links}{The positive analysis transformations}
\begin{proof}
Use the common figure-place notation \(P,M,S\) of the reader's formation
table: \(P\) is the intended conclusion predicate, \(S\) its subject,
and \(M\) the middle. This uniform relabeling preserves the source's
predicate-first sentence order. Every conversion below is CR1, CR2, or CR3.
For the target pattern, major/minor names may switch; where that happens
we explicitly restore the original conclusion by simple conversion.

\emph{Celarent to Cesare and Camestres.} Start with \(PeM,MaS\).
Convert \(PeM\) to \(MeP\). The pair \(MeP,MaS\) is Cesare and
gives \(PeS\). Transposing those same premises gives \(MaS,MeP\),
Camestres with major \(S\) and minor \(P\). Its conclusion is \(SeP\),
whose e-converse is \(PeS\). Both preserve the original universal-privative
problem; the second is the destination not explicitly treated in the
ancient passage according to Boger's table~37.

\emph{Darii to Datisi and Disamis.} Start with \(PaM,MiS\).
Convert the second premise to \(SiM\). The resulting \(PaM,SiM\)
is Datisi and gives \(PiS\). Transpose to \(SiM,PaM\), Disamis with
major \(S\) and minor \(P\). It gives \(SiP\), whose i-converse is
\(PiS\). The latter is another destination marked not treated by the source.

\emph{Ferio to Festino and Ferison.} Start with \(PeM,MiS\).
Converting the major gives \(MeP,MiS\), Festino, concluding \(PoS\).
Converting the minor instead gives \(PeM,SiM\), Ferison, with the same
\(PoS\). No o-conversion is used.

\emph{Cesare to Celarent.} Start with \(MeP,MaS\).
Convert the first to \(PeM\); \(PeM,MaS\) is Celarent with conclusion
\(PeS\).

\emph{Camestres to Celarent.} Start with \(MaP,MeS\).
Convert the second to \(SeM\), then transpose to \(SeM,MaP\).
Celarent gives \(SeP\). Convert that conclusion to \(PeS\).

\emph{Festino to Ferio and Ferison.} Start with \(MeP,MiS\).
Convert the first to \(PeM\) to obtain Ferio. Also converting \(MiS\)
to \(SiM\) gives Ferison. Both target conclusions are \(PoS\).

\emph{Darapti to Darii.} Start with \(PaM,SaM\).
Convert \(SaM\) per accidens to \(MiS\). Darii on \(PaM,MiS\)
gives \(PiS\). The original universal affirmation supplies existent \(M\).

\emph{Felapton to Ferio and Festino.} Start with \(PeM,SaM\).
Convert \(SaM\) per accidens to \(MiS\) to get Ferio. Also converting
\(PeM\) to \(MeP\) gives Festino. Both target conclusions are \(PoS\).

\emph{Disamis to Darii.} Start with \(PiM,SaM\).
Convert \(PiM\) to \(MiP\), transpose to \(SaM,MiP\), and apply
the Darii pattern to obtain \(SiP\). Its i-converse is \(PiS\).

\emph{Datisi to Darii.} Start with \(PaM,SiM\).
Convert the second to \(MiS\), obtaining Darii with conclusion \(PiS\).

\emph{Ferison to Ferio and Festino.} Start with \(PeM,SiM\).
Convert the second to \(MiS\), obtaining Ferio. Also convert the first
to \(MeP\), obtaining Festino. Both conclusions are \(PoS\).

There are seventeen links: two each from Celarent, Darii, Ferio, Festino,
Felapton, and Ferison, and one each from Cesare, Camestres, Darapti,
Disamis, and Datisi. Every new premise has been obtained by conversion;
every target conclusion is either already in the starting orientation or
is restored by simple e/i-conversion. No universal strengthening, o-conversion,
or new reductio subdeduction has been used. This proves all positive links
claimed, not the exhaustive absence of other transformations.
\end{proof}
\end{annexitem}
